% ABGELOEST (2026-09-07): dieser Entwurf ist durch docs/chapter_omol25.tex
% ersetzt (Gruppennamen closed-shell / broken-symmetry, Fig 3 mit Panel A+B,
% Leiter-Tabelle). Nur noch als Verlauf aufbewahrt.
%
% Results, zusammengefuehrt: der Appendix-Teil (Figur fig3 und der doppelte
% Unterabschnitt) ist in den Haupttext gewandert. Aenderungen gegenueber
% dem eingereichten Entwurf, alle formal:
%   - \label{fig:forceerror} -> \label{fig:ferr}   (so wird es 3x referenziert)
%   - \label{fig:drift}      -> \label{fig:barrier} (so wird es referenziert)
%   - Barrieren-Unterabschnitt nur einmal, mit Figur UND Tabelle
%   - "between 337 and 4434 meV, and 5 of 18 beyond" -> die fuenf ueber
%     43 meV reichen von 61 bis 4434, nicht von 337
%   - eV,\AA -> eV\,\AA  (zweimal in der Discussion)
% Alle Zahlen gegen results/*.csv geprueft (check_results_section.py plus
% f_err_max-Mediane, Tiefenspanne, Spearman -0.047).

\section{Results}

\begin{figure}[t]
  \centering
  \includegraphics[width=\linewidth]{figs/Fig1.png}
  \caption{Model residual force against DFT residual force, evaluated at the
  same unrelaxed geometry: the transition state each model produced
  (45~reactions $\times$ 3~models; all 135 structures shown).
  Both axes give $\max_i |F_i|$ over all $3N$ Cartesian components.
  The DFT force is evaluated at the OMol25 level of theory on the surface
  selected by the stability analysis (Sec.~\ref{sec:methods}).
  Colours mark the two classes: $\langle S^2 \rangle$ of that single point is
  exactly~0 for 82 of the 135 structures and at least 0.0579 for the
  remaining~53, so the classes separate without a chosen cut-off.
  Crosses mark group medians, model/DFT, stable then unstable ---
  UMA-S 0.030/0.046, 0.042/0.157; UMA-M 0.019/0.039, 0.030/0.115;
  eSEN 0.042/0.074, 0.038/0.129.}
  \label{fig:silent}
\end{figure}

\subsection{The residual force the model reports does not track the force at its transition state}
Figure~\ref{fig:silent} evaluates both force fields at the identical, unrelaxed
geometry --- the transition state each model produced in its own search.
In the model's own residual, the two classes are nearly indistinguishable: the
median is 0.0299~eV\,\AA$^{-1}$ at the RKS-stable and 0.0344~eV\,\AA$^{-1}$ at
the RKS-unstable structures, a factor of 1.15.
The DFT residual at the same points separates them: 0.0507 against
0.1293~eV\,\AA$^{-1}$, a factor of 2.55.
That unstable-group median is more than twice the criterion the searches
stopped at.
The stable group sets the baseline.
No second solution exists there, so its ratio measures the ordinary
disagreement between model and DFT.
The unstable group exceeds this baseline, and the difference is the effect of
the instability.
Restricting the analysis to the 112 searches that met the band criterion
changes the DFT-side separation from 2.55 to 2.45.
The model output contains no signal that marks the affected structures.
What separates the two groups is one DFT single point with stability analysis
at the model geometry.

\begin{figure}[H]
  \centering
  \includegraphics[width=\linewidth]{figs/fig2_force_mae.png}
  \caption{Mean absolute difference between the model and the DFT force,
$\langle |F_i^\mathrm{MLIP} - F_i^\mathrm{DFT}| \rangle_i$ over all $3N$
Cartesian components, at the transition state each model produced.
Both fields are read at the identical, unrelaxed geometry, so no geometry
difference enters.
The DFT force is one gradient per structure at the OMol25 level of theory,
on the surface selected by the stability analysis
(Sec.~\ref{sec:methods}); colours mark the two classes by $\langle S^2 \rangle$ of that single point (Fig.~\ref{fig:silent}).
Faint dots are individual structures, jittered horizontally only; bars are
group medians, shaded bands their 95\,\% bootstrap intervals
(10\,000 resamples); the arrow gives the ratio of the two medians.}
  \label{fig:ferr}
\end{figure}

\subsection{The learned force field itself is wrong where the restricted solution is unstable}
Fig.~\ref{fig:ferr} measures the distance between the two force fields directly: the vectors are subtracted component by component at the same geometry, and a large value means the fields disagree.
This error separates the two classes for every model, stable against unstable: medians of 0.0104 against
0.0303~eV\,\AA$^{-1}$ for UMA-S, 0.0052 against 0.0177 for UMA-M, and 0.0138
against 0.0268 for eSEN --- factors of 2.9, 3.4, and 1.9.
Pooled, the median grows from 0.0095 to 0.0250~eV\,\AA$^{-1}$, and the largest
single component from 0.0353 to 0.1034~eV\,\AA$^{-1}$.
The overlap is small: 4 of 53 unstable structures lie below the stable median,
8 of 82 stable structures above the unstable one.
The failure in Fig.~\ref{fig:silent} is therefore not a convergence artefact: the force field the models learned differs from the ground-state field at these points --- at the identical geometry, before any optimisation enters.
Within the unstable group the breaking depth spans 0.6 to 3986~meV, yet the force error does not follow it (Spearman $\rho = -0.05$ across the 53
structures): what matters is that a second solution exists, not how deep it lies.

\begin{figure}[H]
  \centering
  \includegraphics[width=0.99\linewidth]{figs/fig3.png}
  \caption{Barrier error at fixed geometry (a) and barrier spread from
  geometry alone (b).
  (a)~One dot is one reaction evaluated with one model: the model barrier
  minus the DFT barrier at the same two unrelaxed structures, reactant and
  transition state, so no geometry error enters.
  Solid bars are means, dashed bars medians; they disagree because in every
  unstable group one named reaction misses by more than 0.5~eV and sets the
  mean almost by itself.
  (b)~One dot is one reaction (45 in total): the same DFT calculation at its
  three model transition states, each barrier against the reactant of the
  same model, closed-shell in every case; plotted is $\max-\min$ over the
  three.
  A reaction counts as unstable when at least one of its three transition
  states has $\langle S^2 \rangle > 0$ in its single point
  (Fig.~\ref{fig:silent}).
  The green line marks 1~kcal/mol (43~meV) for orientation; it is not a
  criterion used in this work.
  DFT at the OMol25 level of theory, surface selected by the stability
  analysis (Sec.~\ref{sec:methods}).}
  \label{fig:barrier}
\end{figure}

\subsection{The energy error is small and the geometry error is not}
Figure~\ref{fig:barrier} and Table~\ref{tab:barrier} split the barrier error
into its two sources: the model's energy error and its geometry error.
Panel~(a) freezes the geometry: the model barrier minus the DFT barrier at
the identical reactant and transition state, so only the energies differ.
This error is small in both classes --- medians of 1.7 to 4.1~meV at the
stable and 6.3 to 13.7~meV at the unstable structures, per model.
The means sit far above the medians in the unstable groups (42.2, 148.9,
174.3~meV) because a single reaction per model misses by more than half an
electronvolt and carries the mean almost alone --- rxn8837 at 0.62~eV for
UMA-S, rxn0894 at 2.52 and 2.69~eV for UMA-M and eSEN; 49 of 53 unstable
rows lie within chemical accuracy.
The models mostly know the energies, even where the forces fail.

Panel~(b) turns the question around and varies only the geometry: the same
DFT calculation runs at the three model transition states of each reaction,
against the closed-shell reactant of the same model, and the spread
$\max - \min$ of the barrier is recorded.
What separates the three numbers is the saddle position and nothing else.
At the stable reactions the median spread is 0.33~meV; at the unstable ones
it is 10.71~meV.
Five of the 18 unstable reactions exceed chemical accuracy, from 61~meV up to
4434~meV (rxn8837), against 2 of 27 on the stable side.
The choice among the three model geometries alone shifts the barrier by up to
4.4~eV, and the models' own residual forces carry no signal of which
geometries are affected (Fig.~\ref{fig:silent}).
A DFT single point at the model geometry repairs the energy part, which was
largely sound, and inherits the geometry part, which was not.
The affected transition states are therefore useful as starting points for a
subsequent optimisation, not as final saddle points.

\begin{table}[!t]
\caption{Barrier errors, in meV.
(a)~Per structure: $|\Delta E^\ddagger_\mathrm{MLIP} -
\Delta E^\ddagger_\mathrm{DFT}|$ with both barriers read at the identical
unrelaxed reactant and transition state, so no geometry error enters.
(b)~Per reaction: $\max-\min$ of the DFT barrier over the three model
transition states, each against the reactant of the same model, closed-shell
in every case; a reaction counts as unstable when at least one of its three
transition states has $\langle S^2 \rangle > 0$.
Chemical accuracy taken as 1~kcal/mol (43~meV), for orientation only.
DFT at the OMol25 level of theory, surface selected by the stability
analysis (Sec.~\ref{sec:methods}).}
\label{tab:barrier}
\centering
\begin{minipage}[t]{0.56\linewidth}
\centering
(a) energy error at frozen geometry,\\ per structure\\[2pt]
\begin{tabular}{llrrr}
\toprule
 & & $n$ & median & MAE \\
\midrule
UMA-S & RKS stable   & 27 & 1.7  & 4.4 \\
      & RKS unstable & 18 & 7.5  & 42.2 \\
UMA-M & RKS stable   & 27 & 2.4  & 2.8 \\
      & RKS unstable & 18 & 6.3  & 148.9 \\
eSEN  & RKS stable   & 28 & 4.1  & 7.6 \\
      & RKS unstable & 17 & 13.7 & 174.3 \\
\bottomrule
\end{tabular}
\end{minipage}\hfill
\begin{minipage}[t]{0.42\linewidth}
\centering
(b) barrier spread from geometry\\ alone, per reaction\\[2pt]
\begin{tabular}{lrrr}
\toprule
 & $n$ & median & $>43$~meV \\
\midrule
all three stable      & 27 & 0.33  & 2 \\
at least one unstable & 18 & 10.71 & 5 \\
\bottomrule
\end{tabular}
\end{minipage}
\end{table}

\subsection{The training labels are not stationary on the surface they describe}
The measurements so far characterise what the models produce; this section
characterises what they learned (measurement in
Sec.~\ref{sec:methods-labels}, results in Table~\ref{tab:hinge}).
Table~\ref{tab:hinge} answers two questions per geometry: the distance from
a saddle of the restricted surface, and from one of the ground-state
surface.
At the stored Transition1x geometries the answer is 0.6~eV\,\AA$^{-1}$ on
both surfaces for the stable reactions --- the two surfaces are the same
there, and 0.6 is what the level change ($\omega$B97M-V/def2-TZVPD labels
on $\omega$B97x/6-31G(d) optimised geometries) alone did to every saddle.
For the unstable reactions the restricted answer is the same 0.6, but on
the ground-state surface, which is unrestricted for these, the answer is
1.64.
How much of the 1.64 is the level and how much is the surface cannot be
told apart at these points, since both changed at once.
The lower block of Table~\ref{tab:hinge} separates them: re-optimising at
the training level restores the saddle on the restricted surface --- the
residual force drops to 0.04 in both groups --- while the ground-state
surface of the unstable group still carries 1.87, a ratio of 32.

\section{Discussion}

Relabelling the inherited geometries carries two potential inconsistencies.
First, the paths are not relaxed at the level the labels are computed at:
every training point sits 0.6~eV\,\AA$^{-1}$ off stationarity, stable and
unstable alike.
Second, for reactions whose restricted solution is unstable, the paths are
not relaxed on the surface the labels describe: Transition1x was computed
restricted, and the saddle of the ground-state surface lies elsewhere.
The measurements separate the two.
The first is not the problem: the stable reactions have the same offset, and
the models learn them to 0.01~eV\,\AA$^{-1}$ (Fig.~\ref{fig:ferr}).
The second is consistent with everything measured: it exists exactly where
the failures concentrate, and it means the models never trained on a
stationary point of the surface they are asked to find saddles on.
One caveat remains: OMol25 contains reactive organic configurations beyond
the recomputed Transition1x paths (RGD1 path frames, MechDB reaction steps,
bond-breaking curves), and we have not tested how close any of them come to
these saddle regions.
None of them, however, was optimised on the surface its labels describe: all
reactive geometries enter as interpolated, force-field- or MLIP-generated, or
inherited frames with unrestricted single-point labels (OMol25, Secs.~G--H).
Whether this is the cause cannot be decided from our data.
It can be tested: re-relax the affected paths on the ground-state surface
and train on them.
